\section{Implementation}

This section describes the software architecture and implementation details of the experimental framework used to validate the proposed approaches. The system was developed in C\# using the .NET runtime environment and designed for high-performance parallel execution.

\subsection{Software Architecture}

The implementation follows object-oriented design principles and consists of several core components organized into a modular architecture.

\subsubsection{Core Data Structures}

The implementation defines several key classes that encapsulate problem entities:

\begin{itemize}
    \item \textbf{Location:} Represents a customer location or depot, containing attributes such as coordinates, demand, service time, and time window $[a_i, b_i]$. Each location is identified by a unique integer ID.
    
    \item \textbf{Vehicle:} Encapsulates vehicle properties including capacity $Q_v$ and operational characteristics. The fleet is represented as a list of Vehicle objects.
    
    \item \textbf{Route:} Represents a vehicle route as a sequence of Location objects, along with computed metrics including current load, total cost, penalties, and vehicle operation time. Each route maintains references to start time and capacity constraints.
    
    \item \textbf{Solution:} Aggregates a collection of Route objects representing a complete assignment of customers to vehicles. The Solution class maintains global metrics: total travel cost, total penalty, and total vehicle operation time. It also stores a reference to the initial greedy solution metrics for comparison purposes.
    
    \item \textbf{Instance:} Encapsulates a complete problem instance including the list of locations, distance matrix, vehicle fleet, and filename. The class provides parsing functionality for Solomon benchmark format files and supports randomized perturbations for stochastic scenario generation.
    
    \item \textbf{Scenario:} Wraps an Instance with a unique scenario identifier, enabling systematic tracking and analysis across multiple stochastic realizations.
\end{itemize}

\subsubsection{Algorithm Implementation}

The solution methods are implemented as separate modules:

\begin{itemize}
    \item \textbf{GreedyApproaches:} Implements the constructive greedy heuristic. The algorithm maintains a single global tour representation (GTR) and iteratively inserts the nearest feasible customer based on distance, time windows, and capacity constraints. When a vehicle reaches capacity or time limits, a new vehicle is dispatched from the depot. The method \texttt{generateGreedySolution(Instance)} returns a complete Solution object.
    
    \item \textbf{TabuSearch:} Implements the metaheuristic optimization procedure. The main method \texttt{run(MaxIterations, TabuSize, Instance, mutationType)} accepts configuration parameters and returns the best solution found. Key features include:
    \begin{itemize}
        \item Initialization with the greedy solution as a starting point
        \item Neighborhood generation using the specified mutation operator (swap, insert, or invert)
        \item Tabu list management using a queue structure with fixed size
        \item Aspiration criterion allowing tabu moves that improve the best-known solution
        \item Early termination after a specified number of non-improving iterations
    \end{itemize}
    
    \item \textbf{NeighborhoodGenerator:} Provides methods for generating neighboring solutions through local search operators. The \texttt{GenerateAllSwaps} method supports three mutation types:
    \begin{itemize}
        \item \textit{swap:} Exchanges two customers between routes or within the same route
        \item \textit{insert:} Removes a customer and reinserts it at a different position
        \item \textit{invert:} Reverses a subsequence of customers within a route
    \end{itemize}
    Each generated neighbor is evaluated immediately, computing cost, penalty, and feasibility metrics.
\end{itemize}

\subsubsection{Scenario Generation and Experimental Framework}

The experimental infrastructure supports large-scale stochastic analysis:

\begin{itemize}
    \item \textbf{InstanceGenerator:} Creates stochastic scenarios by applying random perturbations to customer demands and time windows. The method \texttt{GenerateManyScenarios(numberScenarios, fileNames)} generates a specified number of scenarios for each input file, producing a complete scenario set for statistical analysis.
    
    \item \textbf{DemandSampler:} Provides sampling methods for generating demand variations based on specified probability distributions (normal, uniform, etc.).
    
    \item \textbf{ExperimentRunner:} Orchestrates the execution of experiments across multiple scenarios, parameter configurations, and algorithm variants. The method \texttt{RunExperiments(scenarios, iterations, tabuSizes, mutationTypes, repeats, baseSeed, parallel)} supports:
    \begin{itemize}
        \item Parallel execution using the Task Parallel Library for efficient multi-core utilization
        \item Deterministic reproducibility through explicit random seed management
        \item Comprehensive result collection including objective values, cost components, route structures (GTR), and execution times
    \end{itemize}
    
    \item \textbf{ResultsAnalysis:} Implements statistical analysis and visualization capabilities, including computation of descriptive statistics, percentiles, and empirical cumulative distribution functions (ECDF).
    
    \item \textbf{ScenarioAnalyzer:} Provides clustering and feature extraction functionality using K-means clustering to group scenarios based on structural characteristics (e.g., customer spatial distribution, time window tightness, demand variability).
\end{itemize}

\subsection{Performance Optimization}

Several design choices enhance computational efficiency:

\begin{itemize}
    \item \textbf{Parallel Execution:} The framework leverages the Task Parallel Library to distribute scenario evaluations across available CPU cores. Each scenario is processed independently, enabling linear speedup with the number of processor cores.
    
    \item \textbf{Incremental Evaluation:} Route metrics (cost, penalties, feasibility) are computed incrementally during neighborhood exploration, avoiding redundant full-route re-evaluations.
    
    \item \textbf{Lazy Neighborhood Generation:} The Tabu Search implementation generates neighbors on-demand using LINQ's lazy evaluation (\texttt{.Take(TabuSize*10)}), terminating generation once sufficient candidates are found. The multiplier of 10 provides a reasonable exploration breadth while maintaining computational efficiency.
    
    \item \textbf{Efficient Data Structures:} Tabu list is implemented as a fixed-size queue with O(1) insertion and O(n) membership testing. Routes are represented as lists enabling O(1) indexed access.
\end{itemize}

\subsection{Configuration and Execution}

The main program (\texttt{Program.cs}) defines experimental parameters:
\begin{itemize}
    \item iterations = 100,
    \item tabu size = 30,
    \item number of scenarios = 500,
    \item mutation type - swap.
\end{itemize}

The framework generates 500 stochastic scenarios for each of six Solomon benchmark instances (C101, C201, R101, R201, RC101, RC201), resulting in 3,000 total scenarios. Both greedy and Tabu Search are executed on each scenario with the specified parameters. Results are appended to a CSV file (\texttt{results\_raw.csv}) containing columns for scenario ID, algorithm variant, objective components, route structure, and execution time.

\subsection{Reproducibility and Extensibility}

The implementation ensures reproducibility through:
\begin{itemize}
    \item Explicit random seed initialization (\texttt{baseSeed = 42}, a commonly-used fixed value in computational research for deterministic reproducibility)
    \item Deterministic algorithm execution (fixed iteration order)
    \item Complete parameter logging in output files
\end{itemize}

The modular architecture facilitates extension with additional:
\begin{itemize}
    \item Solution methods (genetic algorithms, ant colony optimization, etc.)
    \item Mutation operators for neighborhood generation
    \item Objective function components (e.g., environmental costs, driver preferences)
    \item Constraint types (e.g., driver breaks, multi-trip routing)
\end{itemize}

The complete source code and experimental data are available in the project repository, enabling full replication of the reported results and supporting future research extensions.